\chapter{Troubleshooting and Common Footguns}
\label{appendix:troubleshooting}

\epigraph{``Experience is simply the name we give our past build failures.''}{--- Oscar Wilde (adapted)}

\noindent
Even the most seasoned functional alchemists occasionally run into a cryptic error or a failed derivation. Here is a curated diagnostic guide for the most common Guix pitfalls.

\section{Footgun 1: ``No such file or directory'' on Existing Scripts}

\begin{tcolorbox}[colback=footgunred!5!white,colframe=footgunred,title=Symptom]
You run an executable Python or Bash script with \texttt{./script.sh} and get:
\begin{lstlisting}[style=guixcli]
bash: ./script.sh: No such file or directory
\end{lstlisting}
even though \texttt{ls -l script.sh} shows the file exists and has executable (\texttt{+x}) permissions!
\end{tcolorbox}

\textbf{Diagnosis}: The script begins with \texttt{\#!/bin/bash}, \texttt{\#!/usr/bin/env}, or \texttt{\#!/usr/bin/python}, but that interpreter does not exist at that absolute path on a pure Guix system.

\textbf{Solution}:
\begin{enumerate}
  \item In a package definition, let \texttt{patch-shebangs} run, or use \texttt{search-input-file}:
  \begin{lstlisting}[style=guixscheme]
  (substitute* "script.sh"
    (("/bin/bash") (search-input-file inputs "/bin/bash")))
  \end{lstlisting}
  \item In a development shell, run the script explicitly with the interpreter:
  \begin{lstlisting}[style=guixcli]
  guix shell python -- python3 ./script.py
  \end{lstlisting}
\end{enumerate}

\section{Footgun 2: Scheme Error: ``Unbound Variable''}

\begin{tcolorbox}[colback=footgunred!5!white,colframe=footgunred,title=Symptom]
When running \guixcmd{guix build} or \guixcmd{guix system reconfigure}, you see:
\begin{lstlisting}[style=guixcli]
error: openssl: unbound variable
hint: Did you forget `(use-modules (gnu packages tls))'?
\end{lstlisting}
\end{tcolorbox}

\textbf{Diagnosis}: Guix organizes packages into granular Scheme modules. Unlike Python where global namespace pollution is common, Guile requires you to explicitly import the module that defines the variable.

\textbf{Solution}: Add the corresponding module to your \texttt{use-modules} or \texttt{use-package-modules} header at the top of the file.

\section{Footgun 3: Test Suite Fails Because of No Network Access}

\begin{tcolorbox}[colback=footgunred!5!white,colframe=footgunred,title=Symptom]
A package build fails during the \texttt{check} phase with socket or DNS resolution errors (e.g. \texttt{getaddrinfo failed}).
\end{tcolorbox}

\textbf{Diagnosis}: The \texttt{guix-daemon} executes builds in a network-isolated container to ensure hermetic reproducibility. Upstream test suites that attempt to ping \texttt{google.com} or download test fixtures on the fly will fail.

\textbf{Solution}:
\begin{enumerate}
  \item If the tests are non-essential network checks, disable the test phase:
  \begin{lstlisting}[style=guixscheme]
  (arguments (list #:tests? #f))
  \end{lstlisting}
  \item Or use \texttt{modify-phases} to disable only the specific offending unit tests:
  \begin{lstlisting}[style=guixscheme]
  (add-before 'check 'disable-network-tests
    (lambda _
      (setenv "SKIP_NETWORK_TESTS" "1")))
  \end{lstlisting}
\end{enumerate}

\section{Footgun 4: Disk Space Running Out in \texttt{/tmp}}

\begin{tcolorbox}[colback=footgunred!5!white,colframe=footgunred,title=Symptom]
Compiling a massive package (like \texttt{gcc}, \texttt{llvm}, or \texttt{chromium}) fails with \texttt{No space left on device} while building in \texttt{/tmp/guix-build-...}.
\end{tcolorbox}

\textbf{Diagnosis}: The build daemon extracts source trees and object files into \texttt{TMPDIR}, which is often mounted on a RAM disk (\texttt{tmpfs}) with limited size.

\textbf{Solution}: Set \texttt{TMPDIR} to a disk directory with ample storage before running the build:
\begin{lstlisting}[style=guixcli]
export TMPDIR=/var/tmp
guix build large-package
\end{lstlisting}

\section{Footgun 5: ``Why didn't \texttt{guix gc} free any space?''}

\begin{tcolorbox}[colback=footgunred!5!white,colframe=footgunred,title=Symptom]
You run \guixcmd{guix gc}, but the disk usage of \texttt{/gnu/store} barely changes.
\end{tcolorbox}

\textbf{Diagnosis}: Old generations of your user profile, system profile, or custom profiles are still referencing those packages as active GC roots.

\textbf{Solution}: First delete old profile generations:
\begin{lstlisting}[style=guixcli]
# Delete old user profile generations
guix package --delete-generations

# Delete old system generations (if using Guix System)
sudo guix system delete-generations

# Reclaim store space
guix gc
\end{lstlisting}
